\documentclass[10pt]{article}
\usepackage[margin=0.84in]{geometry}
\usepackage[T1]{fontenc}
\usepackage{lmodern}
\usepackage{amsmath,amssymb,amsthm,booktabs,microtype}
\usepackage[hidelinks]{hyperref}
\newtheorem{theorem}{Theorem}
\newtheorem{proposition}{Proposition}
\newcommand{\Z}{\mathbb Z}
\newcommand{\Q}{\mathbb Q}
\newcommand{\W}{W(E_6)}

\title{A Dyadic Mark-Saturation Jump for a Primitive Burnside Relation}
\author{Anonymous}
\date{}

\begin{document}
\maketitle

\begin{abstract}
We compute the integral saturation defect created when the primitive
four-versus-four character relation from C63 is measured by the exact C64
table-of-marks map.  The old control directions have Smith invariants
$(2,8)$, while all three character-kernel directions have invariants
$(2,2,8)$.  Their saturation indices are therefore 16 and 32.  A precise
relative quotient, avoiding a subtraction of finite groups of different
ranks, is cyclic of order two and is generated by the normalized
four-versus-four direction.  Every value is recomputed from canonical C63/C64
evidence and checked by two exact Smith-form routes.  The statement is
restricted to the frozen 16-type support.
\end{abstract}

\section{Set-up}
Let $G=\W$ be the order-$51840$ group in the C61 source package, and let
$X_1,\ldots,X_{16}$ be the transitive sets in the C62 dictionary.  C63 gives
the permutation-character matrix and its integer kernel basis
\[
 z_1=X_{10}-X_9,\quad
 z_2=-X_2-X_3-X_5-X_6+X_{11}+X_{12}+X_{13}+X_{14},\quad
 z_3=X_{16}-X_{15}.
\]
The exterior relation is $R_4=-z_2$.  C64 supplies the integral mark map
\(m:\Z^{16}\to\Z^{16}\) on this same support.  Define
\[
 L_o=m\langle z_1,z_3\rangle,\qquad
 L_a=m\langle z_1,z_2,z_3\rangle,
\]
and, for any lattice $L$,
\[
 U(L)=(\mathbb QL)\cap\mathbb Z^{16}.
\]
The scope firewall is \texttt{NO\_BAD\_EULER\_OR\_ROOT\_NUMBER}; no
arithmetic or full-Burnside assertion is made.

\section{Integral method}
For an integer $n\times r$ matrix $V$ of rank $r$, let $d_k(V)$ be the gcd of
all $k\times k$ minors.  The Smith invariants are
\[
 d_1(V),\quad d_2(V)/d_1(V),\quad\ldots,\quad d_r(V)/d_{r-1}(V).
\]
The index of the column lattice in its rational saturation is $d_r(V)$;
these are standard determinantal-divisor facts for integral matrices
\cite{newman1972integral,cohen1993computational}.  Since the C63 character
matrix has rank 13, its kernel is saturated as a subgroup of $\Z^{16}$; a
direct gcd-one maximal minor of the displayed three vectors gives the same
conclusion for the chosen basis.

The producer multiplies the C64 mark matrix by the two kernel submatrices and
enumerates every minor needed for ranks two and three.  No floating point
arithmetic is used.  A separate checker repeats the arithmetic from the
source evidence, and a SymPy exact-SNF backend provides a library-level
cross-check.

\section{Main result}
\begin{theorem}[Dyadic saturation jump]
With $L_o,L_a$ as above,
\[
 \operatorname{SNF}(L_o)=(2,8),\qquad
 \operatorname{SNF}(L_a)=(2,2,8).
\]
Consequently
\[
 U(L_o)/L_o\simeq\mathbb Z/2\oplus\mathbb Z/8,
 \qquad
 U(L_a)/L_a\simeq(\mathbb Z/2)^2\oplus\mathbb Z/8.
\]
\end{theorem}

The exact mark vectors, in the order of the 16 subgroup types, are
\[
\begin{aligned}
m(z_1)&=(0,0,0,0,0,0,0,0,-16,8,0,0,0,0,0,0)^T,\\
m(z_2)&=(0,-4,-2,0,0,2,0,-20,0,0,2,4,-2,0,0,0)^T,\\
m(z_3)&=(0,0,0,0,0,0,0,-4,0,0,0,0,0,0,-2,2)^T.
\end{aligned}
\]
Their contents are 8, 2, and 2.  The normalized vectors
\[
 u_1=m(z_1)/8,\quad u_2=m(z_2)/2,\quad u_3=m(z_3)/2
\]
have gcd-one maximal minors, so they form a saturated basis of $U(L_a)$;
$u_1,u_3$ similarly form a saturated basis of $U(L_o)$.

\begin{proposition}[Relative new class]
There is a canonical finite relative quotient
\[
 U(L_a)/(L_a+U(L_o))\simeq\mathbb Z/2,
\]
generated by the class of $u_2=m(z_2)/2=-m(R_4)/2$.
\end{proposition}

\begin{proof}
Use the saturated basis $(u_1,u_2,u_3)$ of $U(L_a)$.  In these coordinates,
\[
 L_a=\langle(8,0,0),(0,2,0),(0,0,2)\rangle,\qquad
 U(L_o)=\langle(1,0,0),(0,0,1)\rangle.
\]
Therefore $L_a+U(L_o)$ is generated by
$(1,0,0),(0,2,0),(0,0,1)$.  The quotient is the second coordinate modulo
2, so the class of $u_2$ is nonzero and has order two.
\end{proof}

\section{Reproducibility and controls}
The C65 evidence binds C63 hash
\texttt{38f439cfe6ed71616a7c74d68bd07da73f5680566ae16f8c557ab2b5d1d16e26}
and C64 evidence hash
\texttt{7c4673e46f2b97ac03d4e331c762a47286058c36ea243fb20fc39543dd699212},
as well as the C64 scoped manifest.  The producer and checker agree on both
Smith forms; the clean replay returns \texttt{REPLAY\_PASS}; the independent
library route returns \texttt{SNF\_CROSSCHECK\_PASS}; and ten hostile
mutations are rejected.  The exact predecessor bytes, rather than selection
pilots, are the authority.

\section{Boundary}
The result concerns only the 16-type C63/C64 submodule.  It is not a claim
about the full Burnside ring of $G$, an arithmetic field lattice, local
decomposition, Euler factors, root numbers, automorphy, or a Hilbert--Polya
operator.  The one new class is an integral mark-saturation distinction, not
an arithmetic equivalence or inequivalence theorem.

\section{Conclusion}
The C63 character-kernel direction adds exactly one independent dyadic class
to the mark-saturation defect.  The relative quotient formulation isolates
that new class without conflating lattices of different ranks and supplies a
replayable integral invariant for the next adaptive step.

\bibliographystyle{plain}
\bibliography{references}
\end{document}
